\chapter*{Streszczenie w języku polskim}

Celem projektu było wykonanie aplikacji desktopowej do zarządzania hotelem. Program ma pomagać pracownikowi recepcji w kontroli podstawowych danych, takich jak goście, pokoje, typy pokoi, rezerwacje, płatności, pracownicy oraz usługi dodatkowe.

W projekcie zastosowano język C\#, technologię WPF do wykonania interfejsu użytkownika oraz Entity Framework Core do obsługi lokalnej bazy danych SQLite. Aplikacja posiada panel główny ze statystykami, zakładki nawigacyjne, formularze do dodawania danych, tabele z rekordami oraz podstawową walidację pól.

Główne funkcjonalności programu obejmują dodawanie gości, tworzenie rezerwacji, sprawdzanie dostępności pokoi, rejestrowanie płatności oraz dodawanie i edycję usług hotelowych. Końcowym efektem projektu jest działająca aplikacja, która realizuje najważniejsze założenia prostego systemu zarządzania hotelem.
